%=================================%
\section{Linear framework}
\label{sec:linear_method}
%=================================%

In this section we will summarise the linear procedure being used to study eccentric planet-disc interactions. We will begin by reviewing the orbital setup of the embedded planet and also the physical parameters governing the disc profile. We will then outline the linearised equations being solved and introduce the numerical mode finding method. For more details on the setup and an in-depth description of the solution techniques, we refer the reader to Sections 2 and 3 of \paper1 and the references therein.

%---------------------------------%
\subsection{Planet-disc setup}
\label{subsection:planet_disc_setup}
%---------------------------------%

We adopt a single planet with mass $M_\textrm{p}$ orbiting a host star of mass $M_{*}$. The anticipated migration timescale and evolution of orbital elements resulting from the planet-disc interaction are much longer than the orbital timescale, so we assume the planet is on a fixed, eccentric orbit \citep[as per][]{GoldreichTremaine_1980}. This orbit is parameterised by two elements, relative to a coordinate system centred on the star, planetary eccentricity $e$ and the semi-major axis $a$. The Kepler problem then dictates that the orbital mean motion of the planet is $\Omega_\textrm{p} = \sqrt{G M_{*}/a^3}$. The argument of periapse is arbitrarily chosen so it conveniently aligns with the coordinate frame $x$-axis. In this polar reference frame, the planetary radial coordinate is denoted $r_\textrm{p}$, whilst the true anomaly is given by $\phi_\textrm{p}$ (see Fig. 1 in \paper1). Meanwhile, an arbitrary position in the disc is labelled by the radial and azimuthal coordinates $r$ and $\phi$ respectively. 

The gravitational influence of the planet leads to a disturbing potential acting on the disc. In this paper we will consider only the \textit{direct} potential contribution which is given by
%
\begin{equation}
    \label{eqn:direct_potential}
    \Phi(\mathbf{r}) = -\frac{G M_\textrm{p}}{|\mathbf{r}-\mathbf{r}_\textrm{p}|} ,
\end{equation}
%
where $\mathbf{r}$ and $\mathbf{r}_\textrm{p}$ denote the arbitrary position vector relative to the star and the position vector of the planet respectively. Note that we do not account for the \textit{indirect} potential in this work. This contributes an additional term which is proportional to $\cos{\phi}$, owing to the non-inertial reference frame centred on the star rather than the planet-star barycenter. Previous authors often disregard this term \citep[e.g.][]{KorycanskyPollack_1993,PapaloizouLarwood_2000,TanakaEtAl_2002}, so for compatibility when comparing our results, this has also been neglected here. Indeed, \cite{MirandaRafikov_2019a} and \paper1 found that inclusion of the indirect term does not significantly impact the wake morphology.

The second key component of the problem is the gaseous disc with which the embedded protoplanet interacts. We will consider a 2D, coplanar, inviscid disc model which is perhaps the simplest scenario but still contains a richness of physics. The unperturbed, axisymmetric disc is described by smooth power-law profiles in the surface density
%
\begin{equation}
    \label{eqn:surface_density}
    \Sigma(r) = \Sigma_{\textrm{p}}\left(\frac{r}{a}\right)^{-p} ,
\end{equation}
%
and the locally isothermal sound speed 
%
\begin{equation}
    \label{eqn:sound_speed}
    c_\textrm{s} (r) = c_{\textrm{s},\textrm{p}}\left(\frac{r}{a}\right)^{-q/2}.
\end{equation}
%
We adopt a simple locally isothermal equation of state (EoS) which then sets the vertically integrated pressure $P$ according to 
%
\begin{equation}
    \label{eqn:eos}
    P = c_{\textrm{s}}^2 \Sigma ,
\end{equation}
%
such that there is a radial temperature profile
%
\begin{equation}
    \label{eqn:temperature}
    T_\textrm{disc}(r) \propto \frac{P}{\Sigma} = c_\textrm{s}^2 \propto r^{-q},
\end{equation}
%
controlled by the exponent $q$. This choice for the EoS affects the form of the underlying perturbation equations and the resulting wave propagation behaviour as shown by \cite{MirandaRafikov_2020}. According to hydrostatic equilibrium, the aspect ratio $h(r)$ of the disc then follows
%
\begin{equation}
    \label{eqn:aspect_ratio}
    h(r) \equiv \frac{H(r)}{r} = h_\textrm{p} \left(\frac{r}{a}\right)^{(1-q)/2},
\end{equation}
%
where $H(r) = c_s(r)/\Omega_\textrm{K}(r) $ is the pressure scale height and $\Omega_\textrm{K}(r) = \sqrt{GM_{*}/r^3}$ denotes the local Keplerian velocity. Note that all of the quantities subscripted by `$\textrm{p}$' are simply characteristic values evaluated at the planetary semi-major axis $a$. 

Associated with the surface density and temperature structure is an additional pressure support in the radial force balance for the disc. This leads to a slightly sub-Keplerian fluid orbital frequency $\Omega$ due to order $h_{\textrm{p}}^2$ corrections:
%
\begin{equation}
    \label{eqn:fluid_frequency}
    \Omega^2(r) = \Omega_\textrm{K}^2+\frac{1}{r \Sigma}\frac{dP}{dr}=\Omega_\textrm{K}^2\left[1-h_\textrm{p}^2(p+q)\left(\frac{r}{a}\right)^{1-q}\right].
\end{equation}
%
Similarly the fluid epicylic frequency $\kappa$ is slightly detuned from the Keplerian value $\Omega_\textrm{K}$ and is given by
%
\begin{equation}
    \kappa^2 = \frac{2\Omega}{r}\frac{d}{dr}\left(r^2\Omega\right) = \Omega_\textrm{K}^2\left[1-h_\textrm{p}^2(p+q)(2-q)\left(\frac{r}{a}\right)^{(1-q)}\right].
\end{equation}

%---------------------------------%
\subsection{Linear Equations}
\label{subsection:linear_equations}
%---------------------------------%

Linear theory has paved the way for understanding planet-disc interaction since the seminal work of \cite{GoldreichTremaine_1979, GoldreichTremaine_1980}. To be more specific, linearity requires that $M_\textrm{p} \lesssim M_\textrm{th}$ where the thermal mass is defined as $M_\textrm{th} = h_\textrm{p}^3 M_*$. This condition is equivalent to requiring that the Hill radius of the planet is less than the local pressure scale height of the disc.

The mass continuity and Euler equations can then be linearised by decomposing the fluid variables into the steady background and perturbed components in accordance with $X\rightarrow X+\delta X$ (dropping subscript `$0$' from the unperturbed variables). Notably the direct planetary potential acts as a perturbing source term which admits the Fourier expansion
%
\begin{equation}
    \label{eqn:potential_expansion}
    \Phi = \sum_{m = 1}^{\infty} \sum_{l = -\infty}^{\infty} \Phi_{ml}(r)\cos\left(m\phi-l \Omega_\textrm{p} t\right) ,
\end{equation}
%
where only the cosine term appears since we are free to choose a time origin such that the planet is at $\phi = 0$ when $t=0$. The presence of two modal numbers $m$ and $l$ indicates that there are two fundamental periodicities in the problem. Whereas a circular planet appears stationary in a frame corotating with the mean motion, and only requires the azimuthal quantum number $m$, an eccentric planet will execute epicylic oscillations in the frame aligned with the planetary guiding centre and hence require a temporal expansion in $l$ also. This form shows that each $(m,l)$ modal combination contributes an $m$-armed potential forcing with a pattern speed
%
\begin{equation}
    \label{eqn:pattern_speed}
    \omega_{ml} = \frac{l}{m}\Omega_\textrm{p}.
\end{equation}
%
This forcing will drive a corresponding Fourier modal response in the fluid perturbations which are also expanded according to 
%
\begin{equation}
    \label{eqn:modal_repsonse}
    \delta X = \mathrm{Re}\left[\delta X(r)e^{\mathrm{i}(m\phi-l\Omega_\textrm{p}t)} \right],
\end{equation}
%
where for brevity we have dropped subscripts $(m,l)$ from $\delta X_{ml}$. Inserting this ansatz into the linear equations reduces the partial differential equations to a set of coupled ordinary differential equations governing the radial mode structure:
%
\begin{align}
\label{eq:density_pert_eqn}
&    -i \Tilde{\omega}\delta \Sigma + \frac{1}{r}\frac{d}{d r}(r\Sigma \delta u_r)+\frac{i m \Sigma}{r}\delta u_\phi = 0 , 
    \\
\label{eq:ur_eqn}
&    -i \Tilde{\omega}\delta u_r - 2\Omega \delta u_{\phi} = -\frac{1}{\Sigma}\frac{d}{d r} \delta P +\frac{1}{\Sigma^2}\frac{dP}{dr}\delta \Sigma-\frac{d}{d r}\Phi_{ml},
    \\
\label{eq:uphi_eqn}
& -i \Tilde{\omega}\delta u_\phi + \frac{\kappa^2}{2\Omega}\delta u_r = -\frac{i m}{r}\left( \frac{\delta P}{\Sigma}+\Phi_{ml}\right).
\end{align}
%
Here $\tilde{\omega} = m(\omega_{ml}-\Omega)$ is the Doppler-shifted forcing frequency experienced by a fluid element in the disc as the perturbing potential mode sweeps by. These equations can be combined along with the EoS \eqref{eqn:eos} to yield the master equation
%
\begin{align}
\label{eq:master_eqn}
     \frac{d^2}{dr^2}\delta h & +\left\{ \frac{d}{dr}\ln\left(\frac{r\Sigma}{D}\right)-\frac{1}{L_T} \right\} \frac{d}{dr}\delta h \nonumber \\
    & - \left\{ \frac{2m\Omega}{r\Tilde{\omega}}\left[\frac{1}{L_T}+\frac{d}{dr}\ln\left(\frac{\Sigma \Omega}{D}\right)\right] \right. \nonumber \\
    & \left.+\frac{1}{L_T}\frac{d}{dr}\ln\left(\frac{r\Sigma}{L_T D}\right)+\frac{m^2}{r^2}+\frac{D}{c_{\textrm{s}^2}}\right\} \delta h = \nonumber \\
    & = -\frac{d^2 \Phi_{ml}}{dr^2} - \left[ \frac{d}{dr}\ln\left(\frac{r\Sigma}{D}\right)\right]\frac{d\Phi_{ml}}{dr} \nonumber \\
    & + \left\{\frac{2m\Omega}{r\Tilde{\omega}}\left[\frac{d}{dr}\ln\left(\frac{\Sigma \Omega}{D}\right)\right]+\frac{m^2}{r^2}\right\}\Phi_{ml},
\end{align}
%
in terms of the pseudo-enthalpy variable $\delta h = \delta P/\Sigma$. Here $L_\textrm{T}$ denotes the characteristic locally isothermal length scale set by
%
\begin{equation}
    \label{eqn:thermal_scale}
    \frac{1}{L_\textrm{T}} = \frac{d \ln c_\textrm{s}^2}{dr} ,
\end{equation}
%
whilst $D$ measures the offset from a resonance given by
%
\begin{equation}
    \label{eqn:resonance_offset}
    D = \kappa^2-\tilde{\omega}^2.
\end{equation}
%
This formulation of the problem looks equivalent to equations (17)-(21) presented by \cite{MirandaRafikov_2020}, but now $\tilde{\omega}$ and $D$ crucially involve the different pattern speeds $\omega_{ml}$ corresponding to the eccentric planetary motion. When later computing the torque and angular momentum flux, we will use the original fluid variables which are given by  
%
\begin{align}
    \label{eq:dsigma_dh}
    \delta\Sigma & = \frac{\Sigma \delta h}{c_{\textrm{s}^2}} , \\
    \label{eq:dur_dh}
    \delta u_r &= \frac{i}{D} \left[ \left( \Tilde{\omega}\frac{d}{d r}-\frac{2m\Omega}{r}\right)(\delta h+\Phi_{ml})-\frac{\Tilde{\omega}}{L_T}\delta h\right] , \\
    \label{eq:duphi_dh}
    \delta u_\phi & = \frac{1}{D} \left[ \left( \frac{\kappa^2}{2\Omega}\frac{d}{d r}-\frac{m\tilde{\omega}}{r}\right)(\delta h+\Phi_{ml})-\frac{\kappa^2}{2\Omega L_T}\delta h \right].
\end{align}
%

%---------------------------------%
\subsection{Review of resonances}
\label{subsection:resonances}
%---------------------------------%

Various singularities are present in equation \eqref{eq:master_eqn} at the radii where $\omega_{ml}=\Omega$ (i.e. $\tilde\omega=0$) or $D=0$, which evidences the existence of resonant locations. Physically they correspond to a maximal coupling between the forcing due to the perturbing potential and the free oscillations supported by the disc -- thereby yielding sites where waves are launched and angular momentum is efficiently transferred between the planet and fluid \citep[e.g.][]{GoldreichTremaine_1979}. These resonances come in a variety of flavours. 

%.................................%
\subsubsection{Lindblad resonances}
%.................................%

Nominal Lindblad resonances occur when $D = 0$ such that $\Tilde{\omega}=m(\omega_{ml}-\Omega) = \pm \kappa$, see equation (\ref{eqn:resonance_offset}). In other words the Doppler shifted forcing frequency matches the fluid epicyclic frequency. In a disc with uniform pressure (i.e. $q=p=0$) in which $\Omega = \kappa =\Omega_\textrm{k}$, this corresponds to the standard mean motion resonant locations
%
\begin{equation}
    \label{eqn:lindblad_resonances}
    \frac{r_{\pm}}{a} = \left(\frac{m\pm 1}{l}\right)^{2/3}.
\end{equation}
%
Lindblad resonances are the locations where the inertial-acoustic waves are launched, which then propagate away from the planet. For suitably low values of $m \ll r/H$, \cite{GoldreichTremaine_1979} used the WKB approximation to describe the wave launching as occurring strictly at these resonances, although \citet{RafikovPetrovich_2012} have shown that this simple picture is incomplete, which manifests itself in the \textit{negative torque density} phenomenon. 

However, for $m \gtrsim r/H$ the highly non-axisymmetric wave structure modifies the underlying dispersion relation wherein the azimuthal pressure gradients become significant. This causes the effective wave launching resonant positions to be displaced from their fiducial Lindblad value and instead pile up at a  distance of $\sim 2H/3$ from the planet \citep{Artymowicz_1993}. This coincides with the location at which the azimuthal disc shear flow becomes supersonic relative to the planet \citep{GoodmanRafikov_2001} and accounts for the \textit{torque cut-off} phenomenon, wherein the torque contribution from the high-$m$ potential harmonics becomes exponentially suppressed. This deviation of the wave launching from the classical \cite{GoldreichTremaine_1979} WKB approach can lead to notable discrepancies with reality when the higher-$m$ modes are important.   

%.................................%
\subsubsection{Corotation resonances}
\label{subsubsec:corotation}
%.................................%

Another important resonant location occurs where $\tilde{\omega} = 0$ such that $\omega_{ml} = \Omega$. In other words the pattern speed matches the local orbital velocity and they move together, lending to the name \textit{corotation} resonances. For a circular planetary orbit, with only one pattern frequency, the corotation resonance coincides with the coorbital radius to $\mathcal{O}(h_\textrm{p}^2)$. However, eccentric orbits are decomposed into Fourier modes exhibiting a variety of pattern speeds, corresponding to a host of non-coorbital corotation resonances \citep{GoldreichTremaine_1980} which also need to be considered. 

Such corotation resonances are complicated locations in the planet-disc interaction problem which are interpreted in a variety of different regimes \citep[e.g.][]{GoldreichTremaine_1979, Ward_1991, PaardekooperPapaloizou_2009}. Since the methodology in this paper is an entirely linear analysis, we are limited to probing the linear corotation torque as originally investigated by \cite{GoldreichTremaine_1979}. In this regime, the corotation resonance manifests itself via a discontinuity in angular momentum flux across the resonant location (see Section \ref{subsection:circular_wake_amf}), equal to the torque imposed in the vicinity of corotation. The angular momentum flux excited in this region decays evanescently from corotation, indicating that the angular momentum is deposited locally and cannot escape from the vicinity of the resonance; this is in stark contrast to the wavelike propagation of spiral waves away from Lindblad resonances. If efficient communication of this deposited angular momentum is maintained through viscous or diffusive effects, the linear corotation torque \citep{PaardekooperEtAl_2011} is preserved. Otherwise, nonlinearities will develop and manifest as the \textit{horseshoe drag} wherein advective momentum terms govern the exchange of angular momentum as fluid parcels librate around corotation \citep[e.g.][]{Ward_1991,PaardekooperEtAl_2010}. Whilst we will not resolve this nonlinear dynamics in this study, we find that the linear corotation torque is often subdominant compared with the Lindblad torque, see Section \ref{sec:eccentric_torques}. Furthermore, we exploit a standard torque decomposition technique to separate the Lindblad and corotation contributions \citep[see e.g.][]{MirandaRafikov_2020}, which can then be dealt with separately if desired.

%---------------------------------%
\subsection{Numerical methodology}
\label{subsection:numerical}
%---------------------------------%

%.........................%
\subsubsection{Summary}
%.........................%

In order to compute the net wake structure, and associated torque back-reaction onto the planet, we solve equation \eqref{eq:master_eqn} numerically. This immediately allows us to circumvent the compromising assumptions entailed with analytical approaches. In particular, the classic torque formula of \cite{GoldreichTremaine_1979} breaks down when azimuthal pressure gradients enter the problem for large $m\sim h_{\textrm{p}}^{-1}$. By solving for the global structure of the modes directly, we fully capture the excitation of the density wave and sample its interaction with the perturbing potential over the finite excitation interval, rather than just at isolated resonances. Motivated by this we can readily exploit the numerical method tested and used in \paper1. This method implements a solution technique largely inspired by the work of \cite{KorycanskyPollack_1993} and developed for circular planet-disc interactions by \cite{MirandaRafikov_2020}. For full details we refer the reader to these works and only summarise the key steps here. 

Firstly, the orbit of the planet is specified by a numerical solution to Kepler's equation. Since we are interested in probing larger values of the eccentricity, we cannot simply use the leading order epicylic approximation which expands the potential perturbation in terms of Laplace coefficients \citep[e.g.][]{GoldreichTremaine_1980}. Indeed, as eccentricity increases, this expansion converges slowly and becomes inaccurate. Instead, the perturbing potential is defined across the disc at each time during the planetary orbit through equation \eqref{eqn:direct_potential}. By direct numerical integration across the temporal and azimuthal periods, we numerically extract the Fourier coefficients $\Phi_{ml}(r)$ at each radius in the disc according to
%
\begin{align}
    \label{eq:phi_mode_integral}
    \Phi_{ml}(r) &= \frac{1}{2\pi^2}\int_{t = 0}^{2\pi/\Omega_\textrm{p}} \int_{\phi = 0}^{2\pi} \Phi(r,\phi,t) \cos(m\phi-l \Omega_\textrm{p} t) \,d\phi\,dt \nonumber\\
             &= -\frac{GM_{\text{p}}}{2\pi^2}\int_{t = 0}^{2\pi/\Omega_\textrm{p}} \int_{\phi = 0}^{2\pi} \frac{\cos(m\phi-l \Omega_\textrm{p} t) \,d\phi\,dt}{\left[r^2+r_{\text{p}}^2-2 r r_{\text{p}} \cos\psi+\epsilon^2\right]^{1/2}} . 
\end{align}
%
Here $\psi(t) = \phi-\phi_{\text{p}}(t)$ denotes the azimuthal displacement from the planet and $\epsilon(|{\bf r}-{\bf r}_\textrm{p}|)$ is the softening kernel, employed to prevent singularities at the location of the planet (when $r - r_{\text{p}} = \psi = 0$) and to mimic the vertical averaging of the three dimensional disc in this 2D setup. Throughout this work we will use the simplest Plummer softening prescription $\epsilon = b H(r_\textrm{p})$, for a choice of constant parameter $b$. Recent work usually assumes the softening to be a significant fraction of the local scale height $H_\textrm{p}$, i.e. $b\sim \mathcal{O}(1)$, so as to simulate the effects of a 3D potential \citep[e.g.][]{PaardekooperPapaloizou_2009}. 

This method for numerical extraction of potential harmonics $\Phi_{ml}$ has proven successful in our previous synthesis of wake morphologies in \paper1. Note that in equation \eqref{eq:phi_mode_integral} we are neglecting the indirect potential component but do compute the direct $m=1$ term. Inspection of equation \eqref{eqn:lindblad_resonances} shows that the nominal inner Lindblad resonance vanishes for this Fourier component, which makes the inner boundary condition numerically unstable and a higher order scheme must be used to solve for the wave. Whilst this $m=1$ term does not change the qualitative wake morphology, we find that it does provide a non-negligible contribution to the torque quantities and should be included (see Fig.~\ref{fig:circular_torques} and discussion in Section \ref{subsection:circular_theory}).

The potential harmonics $\Phi_{ml}$ are evaluated over a discrete radial grid and then fed into the master equation \eqref{eq:master_eqn} for a given choice of $m$ and $l$. The potential at arbitrary values of $r$ can then be interpolated as required. A pair of homogeneous, unforced solutions as well as the forced, inhomogeneous response are integrated outwards from the mode-specific corotation radius where $\omega_{ml}=\Omega$ towards the inner and outer boundaries at $r/a = 0.05$ and $5.0$ respectively. The relative contributions from the unforced wave components are fixed by specifying outgoing radiative boundary conditions at the inner and outer disc edges. Spurious incoming waves are removed by adjusting the boundary conditions using a phase-gradient error minimisation technique \citep{KorycanskyPollack_1993}. The net wake is finally constructed as a superposition of all these individual modes.

%.........................%
\subsubsection{Convergence}
\label{subsubsection:convergence}
%.........................%

Qualitatively, integers $m$ and $l$ govern the azimuthal and temporal convergence of our solutions, respectively. For a circular planet, the wake structure is stationary in a frame which corotates with the orbit, and hence $l=m$, i.e. the pattern speed of modes is equal to the mean motion $\Omega_\mathrm{p}$. However, as planetary eccentricity becomes non-zero, epicylic motion of the planet about the guiding centre introduces an essential time-dependence and `off-diagonal' $|l-m|\neq 0$ modes begin to contribute to forcing. Indeed, the classic disturbing Function for a perturbing planet with small $e$ has a modal expansion which goes as $\Phi_{lm}\propto e^{|l-m|}$ \citep[see][]{MurrayDermott_1999}. Therefore, as we boost $e$ to larger values, we need to include modes of higher order in $|l-m|$ in our analysis. Furthermore, whilst the qualitative morphology of the wake structure might appear to converge rather quickly, special care should also be taken to ensure that the second order quantities such as the torque and angular momentum flux have converged as well when $m$ and $|l-m|$ increase. This ever more challenging convergence criterion with increasing eccentricity ultimately limits our ability to probe arbitrarily high values of $e$. We therefore  restrict our attention to cases where we are confident of convergence, as discussed in more detail in Appendix \ref{app:torque_converegence}. Furthermore, during the post-processing of the many thousands of calculated modes, we very occasionally find some (numerically) spurious behaviour in a few $m=1$ waves which would produce slight deviations in the broad trends discussed later in this paper. We discuss this anomalous behaviour in Appendix \ref{app:m1q1}. Ultimately, these modes are filtered out in our final results, as justified by the compatibility of the adjusted results with adjacent points in parameter space.

%.........................%
\subsubsection{Computing the torque and angular momentum flux}
\label{subsubsection:torque_computation}
%.........................%

With the wave modes in hand, we are in a position to construct second order quantities such as the torque density and angular momentum flux (AMF). These are second order in that they involve a product of linear perturbations. Indeed, the AMF is defined as 
%
\begin{equation}
    F_{J}(r)  = r^2 \Sigma(r) \oint \operatorname{Re}[\delta u_r(r,\phi)] \operatorname{Re}[\delta u_\phi (r,\phi)] d\phi ,
    \label{eq:AMF-def}
\end{equation}
%
where the net velocity perturbations are given by a sum over modes
%
\begin{equation}
    \delta u_r(r,\phi) = \sum_{m,l} \delta u_{r,ml}(r) \exp{[\mathrm{i}(m\phi-l\Omega_\textrm{p} t)]},
\end{equation}
\begin{equation}
    \delta u_\phi(r,\phi) = \sum_{m^\prime,l^\prime} \delta u_{\phi,m^\prime l^\prime}(r) \exp{[\mathrm{i}(m^\prime\phi-l^\prime\Omega_\textrm{p} t)]}.
\end{equation}
%
Inserting these expressions into the equation (\ref{eq:AMF-def}) and simplifying we see that 
%
\begin{align}
    F_{J}(r)  = \pi r^2 \Sigma(r) \sum_{m} & \sum_{l,l^\prime} \operatorname{Re}[\delta u_{r, ml} \delta u_{\phi, ml^\prime}^{*}] \cos[(l-l^\prime)\Omega_\textrm{p}t]\nonumber\\
    & +\operatorname{Im}[\delta u_{r, ml} \delta u_{\phi, ml^\prime}^{*}] \sin[(l-l^\prime)\Omega_\textrm{p}t],
\end{align}
%
where $^{*}$ denotes the complex conjugate. Similarly the radial torque density is defined as
%
\begin{equation}
    \frac{dT}{dr} = -r \oint \operatorname{Re}[\delta\Sigma(r,\phi)]\operatorname{Re}\left[\frac{\partial \Phi_\textrm{p}}{\partial\phi}\right]d\phi ,
\end{equation}
%
which can be written in terms of the modal contributions as
%
\begin{align}
\frac{dT}{dr} = -\pi r \sum_{m} \sum_{l,l^\prime} m & \left\{ \cos[(l-l^\prime)\Omega_\textrm{p}t] \Phi_{ml^\prime} \operatorname{Im} [\delta \Sigma_{ml}] \right. 
\nonumber\\
& \left. -\sin[(l-l^\prime)\Omega_\textrm{p}t] \Phi_{ml^\prime} \operatorname{Re} [\delta \Sigma_{ml}]\right\}.
\label{eq:time_dependent_dTdr}
\end{align}
%
Notice that the azimuthally-averaged quantities fix $m=m^\prime$ so only equal azimuthal wavenumbers couple together and contribute to the orbit-averaged torque. Furthermore, the temporal dependence enters through the sinusoidal argument $(l-l^\prime)$ and so the AMF and torque density will oscillate over the orbital period (see eccentric results in Section \ref{subsection:radial_torque_structure} and Fig.~\ref{fig:time_dependent_torque}). We are primarily interested in the cumulative, secular effects over many orbital periods and thus average over the mean motion period giving us
%
\begin{align}
    & \langle F_{J}(r) \rangle = \sum_{m,l} F_{J, ml}, \\
    %
    &\left\langle \frac{dT}{dr} \right\rangle = \sum_{m,l} \frac{dT_{ml}}{dr} ,
\end{align}
%
where $\langle\cdot\rangle = (\Omega_\textrm{p}/2\pi)\oint\cdot dt$ and we define the individual modal contributions to be
%
\begin{align}
    \label{eq:F_jml}
    F_{J, ml} & = \pi r^2 \Sigma(r) \operatorname{Re}[\delta u_{r, ml} \delta u_{\phi, ml}^{*}], \\
    \label{eq:dTdr_jml}
    \frac{dT_{ml}}{dr} &= -\pi r m \Phi_{ml} \operatorname{Im} [\delta \Sigma_{ml}].
\end{align}
%
Henceforth we shall drop the angled brackets and simply remember that we are dealing with the averaged torque and AMF quantities in the remainder of this paper, unless stated otherwise. We also define the time-averaged, radially-integrated torque to be
%
\begin{equation}
    \label{eq:one_sided_torque}
    T(r) = \int_a^{r} \frac{dT}{dr} dr.
\end{equation}
%
Although this one-sided integral differs compared with some previous conventions, it has been used recently by \cite{CimermanEtAl_2024} to clearly partition the torque contributions interior and exterior to the planet. Indeed with this definition we are integrating outwards from the guiding centre radius so that the net torque acting on the entire disc is given by the difference (rather than the sum) $T_\textrm{net} = T(r_\textrm{out})-T(r_\textrm{in})$. Here, $r_\textrm{in}$ and $r_\textrm{out}$ denote the inner and outer disc boundaries respectively. For example, a positive value of $T(r>a) > 0$ arises from positive torque density contributions depositing angular momentum in the outer disc. The equal and opposite back reaction on the perturber then extracts planetary angular momentum and (for a circular orbit) promotes inwards migration. Meanwhile a positive value of $T(r<a) > 0$ owes to a negative torque density which extracts angular momentum from the inner disc. This is accordingly injected into the planet and (for a circular orbit) drives outwards migration. This sign convention is chosen to allow us to easily compare the asymmetry of the inner and outer disc torques (as seen in Fig.~\ref{fig:circular_morphology_torques} and \ref{fig:torque_density_integrated_torque}) -- highlighting the `repulsive' competition either side of the guiding centre radius.

It is also useful to decompose this net torque into the separate Lindblad and corotation contributions using a recipe similar to that employed by \cite{MirandaRafikov_2020}. By linearity, the total torque acting over the full disc can be calculated for each separate modal contribution
%
\begin{equation}
    T_{\textrm{net},ml} = \int_{r_\textrm{in}}^{r_\textrm{out}} \frac{dT_{ml}}{dr}dr = T_{ml}(r_\textrm{out})-T_{ml}(r_\textrm{in}),
\end{equation}
%
which are then summed to yield the net torque $T_\textrm{net} = \sum_{m,l} T_{\textrm{net},ml}$. At the inner and outer boundaries, sufficiently far from the planet, the torque contributions have plateaued (see discussion in Section \ref{subsection:circular_wake_amf} and the black lines in the associated Fig.~\ref{fig:circular_morphology_torques}). 

The corotation torque is now extracted from every $T_{ml}$ modal combination by looking for the AMF discontinuity in the vicinity of each $\omega_{ml} = \Omega$ resonance. Since we only want to extract the discontinuity about a specific point, we must precisely pinpoint the location of the corotation resonance. To second order in the aspect ratio $h_\mathrm{p}$, the corotation resonance occurs at
%
\begin{equation}
    r_{\textrm{C},ml} = a \tilde{r}_{\textrm{C},ml}\left[1-\frac{1}{3} h_\textrm{p}^2 (q+p) \tilde{r}_{\textrm{C},ml}^{1-q}\right],
\end{equation}
%
where $\tilde{r}_{\textrm{C},ml} = (l/m)^{-2/3}$ is the nominal corotation resonance location in a pressureless disc normalised by $a$, for a given $(m,l)$ mode. We then numerically interpolate the calculated $F_{J,ml}(r)$ and $T_{ml}(r)$, which are evaluated at discrete grid locations, onto points on either side of $r_{\textrm{C},ml}$ such that 
%
\begin{align}
    &\Delta F_{J,ml} = F_{J,ml}(r_{\textrm{C},ml}+\Delta r) - F_{J,ml}(r_{\textrm{C},ml}-\Delta r),\\
    &\Delta T_{ml} = T_{ml}(r_{\textrm{C},ml}+\Delta r) - T_{ml}(r_{\textrm{C},ml}-\Delta r),
\end{align}
%
where the narrow width $\Delta r = 0.005a$ is chosen to encompass the AMF discontinuity. The modal corotation torque acting on the disc is then computed via the difference
%
\begin{equation}
    T_{\textrm{C},ml} = \Delta T_{ml}-\Delta F_{J,ml},
\end{equation}
%
where the inclusion of the torque differential here offsets the change in AMF due to $dT_{ml}/dr$ over this small but finite interval. Finally, the remaining modal Lindblad torque is computed as
%
\begin{equation}
    T_{\textrm{L},ml} = T_{\textrm{net},ml}-T_{\textrm{C},ml}.
\end{equation}
%
This decomposition process is repeated separately for each modal combination before the individual results are finally summed to obtain the partitioning of the net torque into the total corotation and Lindblad contributions; $T_\textrm{C} = \sum_{m,l} T_{\textrm{C},ml}$ and $T_{L} = \sum_{m,l} T_{\textrm{L},ml}$, respectively.




